\section{Transporting patched support through a finite chain}

Work with the original normalized patching data
$A:\mathrm{TWAut}(u,J,\mathrm{Unit},D,T)$ of $\mathrm{Core}$.
Write $q=T.q$ and $D_0=D+nq$. There are $h+1$ points and
$h$ connecting primes. Their integral rings are $K_i[[t]]$,
with $O\twoheadrightarrow K_i$ and the prescribed section
sent to $t$. The points may have different block sizes and different
numbers $g_i$ of classes in their auxiliary cohomology data.
We keep the same ideal $J$ throughout.

% BEGIN R70P1
\paragraph{R70P1. Simultaneous lifts with prescribed prime kernels.}
Use LIM18, PRG13 and the specialization maps of PRG23 on this
Taylor--Wiles system. Write $R=W.R$, $I=\mathrm{ii}(W)$,
$C=R/I$ and $S=\mathrm{Spo}(A)$. PMS30--31 give a surjection
$\bar\pi:C\twoheadrightarrow S$ inducing $\pi:R\to S$.
By PRG15--18, $S$ is a finite free local $B$-algebra of positive
rank generated over $B$ by the labeled splitting roots.
Write $p_i=\ker f_i$. Given the original connecting primes
$e_i\subset p_i,p_{i+1}$, take a prime $P_0$ over $p_0$ by
lying over. Given $P_i$, flat going down supplies $E_i\subset P_i$
over $e_i$. Integral going up supplies $P_{i+1}\supset E_i$ over
$p_{i+1}$. At the two ends apply going down once more to lift
the given starting prime and any specified target prime below the
selected endpoint primes.

For each $P_i$ there is a map $f_i':S\to K_i[[t]]$ above $f_i$
with kernel exactly $P_i$. Indeed $S/P_i$ is finite integral over
$B/p_i$, and its fraction field embeds in an algebraic closure
of $K_i((t))$ over $B/p_i$. The specialized monic polynomials
already split into the specified roots in $K_i[[t]]$. The images
of all generators are therefore among those roots. Their $B$-algebra
is in $K_i[[t]]$, as claimed. Equality of the two monic split
polynomials implies that these roots, at each auxiliary index,
are a permutation of the distinct prescribed list. This applies
separately on the predetermined prime $P_i$, without a new
choice of connecting prime.

Use the single quotient pair of SCI2--6: its residual characters
are distinct because the residual pair has distinct characters and
the group projection is surjective. No block size of $x_i$ enters
that multiplicity-free determinant argument. For each $i$, applying
the field-by-field construction of SCI9--10 to $f_i'\bar\pi$
gives an ordinary point. Its local characters and determinant
are those of $x_i$ through the group projection. Its $U$'s are
the preceding root images, its $y$'s map to zero, and its section
maps to $t$. The local and auxiliary trace relations in the
definition of its kernel are the same relations of the single
quotient pair for every $i$. In particular the kernel of the
point map is $\bar\pi^{-1}P_i$ and the root-difference product
avoids $P_i$.

PMS15--16 and PMS29 apply together to all $h+1$ primes.
Write $C_i$ for the completion of $C$ localized at
$\bar\pi^{-1}P_i$, and $S_i\subset C_i$ for its support ideal.
Their common dimension bound on minimal primes of support is
$1+D+nq=1+D_0$. On all primes containing the auxiliary parameters
the support test is contraction of $\operatorname{Ann}_R X_0$.
PMS11 and PMS15 give movement between any two of these centers,
with equality of the contractions of the supported primes
they move. PMS32 identifies the augmentation support over any
prime $P$ of $S$ by
\[
 \operatorname{Ann}_R X_0\subseteq\pi^{-1}P
       \quad\Longleftrightarrow\quad J\subseteq P\cap B .
\]
All the variable images for the auxiliary parameters vanish under
$\pi$. These statements apply for $h=0$ as well.
% END R70P1

% BEGIN R70P2
\paragraph{R70P2. A thick prime remains thick above the limiting relation set.}
Write $G'$ for the limiting index-two group, $G$ for the original
one and $p:G'\twoheadrightarrow G$ for the continuous projection.
Write $I_u$ for the original closed ordinary Cayley--Hamilton
ideal in $B[G]$ and $I_o$ for the closed ideal in $C[G']$ including
the two auxiliary relations and their polarized images. The map
from $C[G']$ to $S[G]$ induced by $\bar\pi$ and $p$ takes
$I_o$ into the ordinary ideal generated by $I_u$ over $S$.

First check the algebraic generators. The determinant on the
limiting group is the pullback of the specialized determinant.
Hence every coefficient of a generic Cayley--Hamilton identity
maps into its Cayley--Hamilton ideal, by BCN16. The ordered
ordinary factors map to their original ordered factors, and
an inertia relation maps to one of the original inertia
relations by adaptedness. Polarization commutes with both maps
on basis elements by cyclotomic and conjugation compatibility;
the polarized ordinary and inertia generators thus map to the
corresponding generators. At index $j$ the first auxiliary
relation maps to the characteristic polynomial of $[T.\mathrm{sg}_j]$
evaluated at that element: its factors are the splitting roots
in $S$. This is a Cayley--Hamilton relation. Its polarized image
also belongs to the closed ordinary ideal. Indeed the pair's polarization identifies the two full laws: on every group word
the two laws have equal characteristic coefficients, and the
integral word identity of JSP9 determines the coefficients on a
generic linear combination. The star anti-involution therefore carries every generic
Cayley--Hamilton relation to one for the transformed element. The second auxiliary relation maps to zero
because $\mathrm{ta}_j$ maps to one and $y_{ji}$ to zero.
Its polarized image is zero as well.

For the closure assertion fix an open normal subgroup $H$ of
$G$ and an adic depth $N$ in $S$. Use $p^{-1}H$ in $G'$.
The map $\bar\pi$ is local and sends $\mathfrak m_C^N$ into
$\mathfrak m_S^N$. The preceding generator check places
the image in the approximation to the original closed ideal
over $S$ at $(N,H)$: each already closed original relation
belongs to this approximation. Thus it lies in their intersection.
The maximal ideal topologies of $B$ and $S$ are cofinal on $S$.
Indeed $S/\mathfrak m_BS$ is an Artinian local ring, so for
some $c$ we have $\mathfrak m_S^c\subseteq\mathfrak m_BS$.
BCN17 applies to the finite free extension $S/B$ and identifies
the closed ideal with $SI_u$ in $S[G]$. This proves the assertion.

Let $P$ be any prime of $S$. Given a prime $Q$ in any commutative
$C$-algebra $H_1$ with the same contraction to $C$ as $P$ through
$\bar\pi$, suppose $P\cap B$ is thick for $I_u$. Choose a field
witness $B\to k_0$ for this nonsemisimplicity. Amalgamate $k_0$
and $\operatorname{Frac}(S/P)$ over $\operatorname{Frac}(B/(P\cap B))$
by taking a prime quotient of their nonzero tensor product and
a fraction field. Amalgamate the result and
$\operatorname{Frac}(H_1/Q)$ over $\operatorname{Frac}(C/(P\cap C))$
in the same way. Both operations preserve the field embeddings.
In the final field $k_*$ the map $H_1\to k_*$ has kernel $Q$.
There the specialized limiting group algebra modulo the image of
$I_o$ surjects onto the specialized original group algebra modulo
$I_u$, by the first part and the surjectivity of $p$.
The latter algebra is the scalar extension of the witness algebra
from $k_0$ to $k_*$. Its semisimplicity would descend by BCN20.
A quotient of a semisimple ring is semisimple, so the limiting
algebra is nonsemisimple too. This proves the native
$\mathrm{nonssT}$ for every completion, with a witness field
at the actual prime $Q$.
% END R70P2

% BEGIN R70P3
\paragraph{R70P3. Moving a supported discriminant prime along every edge.}
On $C$ call a prime discriminant if it avoids a determinant
$\det(\operatorname{tr}([a_i b_j]))$ for two lists of $n^2$
group elements of the limiting group. The condition is inherited
by smaller primes. At a preimage of a prime of $S$ this is
equivalent to the original discriminant condition over $B$:
the determinant laws match and the projection on the groups is
surjective, so one can project or lift both lists.

Suppose on each actual $C_i$ one has the following local comparison.
If $Z,Q$ are primes, the contraction of $Z$ is discriminant and
$\dim(C_i/Z)\ge1+D_0$, then $Z\subseteq Q$ whenever the
contraction of $Q$ is discriminant or $Q$ is thick for the
limiting relation set. This is the incidence implication in the
native $\mathrm{GalOut}$ applied with the three local cohomology
conditions $K1,K2,K3$. It refers to the very completion and
ordinary point of R70P1.

The initial discriminant prime containing $J$ lifts below $P_0$
by R70P1. Pull it back to $C$ and lift further to $C_0$ by
faithful flatness of local completion. The contraction contains
the auxiliary parameters. By PMS32 and the fiber equivalence
of R70P1 its lift contains $S_0$ (here $S_0$ denotes the support
ideal in $C_0$). Thus there is a supported prime with discriminant
contraction at the first vertex.

Suppose such a prime has been obtained in $C_i$. Choose below it
a minimal prime $Z$ over $S_i$. Its contraction is still
discriminant and it satisfies the dimension bound. Lift the edge
prime $\bar\pi^{-1}E_i$ to a prime $Q$ of $C_i$ by the same
flatness. It is thick by R70P2, so the local comparison gives
$Z\subseteq Q$. Its contraction to $R$ is below the next center
since $E_i\subseteq P_{i+1}$, and all parameter images are in
$Q$. PMS11 moves $Z$ to a prime containing $S_{i+1}$ with the
identical contraction to $C$ (the map $R\to C$ is onto).
This supplies the induction hypothesis at $i+1$. In particular
no claim about equality of separately chosen local minimal
primes is needed.

At the last center lift the specified target $q'\subseteq p_h$
through its chosen prime of $S$ and then to $Q$ in $C_h$.
A discriminant target remains discriminant, and a thick target
gives a thick $Q$ by R70P2. Under the last supported prime
choose a minimal support prime as before. The comparison puts it
in $Q$. All parameters lie in $Q$, so the fiber equivalence and
PMS32 give $J\subseteq q'$. This proves the complete finite-chain
conclusion including the single-vertex case, with the original
quantification over all starting, connecting and terminal primes.
BCN23 and PQL20 now give $\mathrm{Raw}$ provided the displayed local
comparison is supplied at all normalized vertices. This conditional
step allows arbitrary $g_i$, $h$ and admissible rank.
% END R70P3

% BEGIN R70P4
\paragraph{R70P4. The positive-dimensional local relation matrix.}
We prove $\mathrm{presOut1}$. Let $O$ be the complete coefficient DVR
of the local presentation, with uniformizer $\varpi$, and let
$B$ be a complete Noetherian local $O$-algebra with the same
residue field. In this paragraph $B$ is the local presentation
ring. Suppose the finite commutative $B$-algebra $A$ contains
$B$ and splits as $B\oplus N$ as a $B$-module, with products of
two elements of $N$ lying in $\mathfrak m_B$. Let $I_N\subset B$
be the ideal generated by these products and put
$J_2=(\varpi)B+\mathfrak m_B^2+I_N$. The presentation data
consist of $D_0$ elements $X_a\in\mathfrak m_B$ with
$\mathfrak m_B\subseteq(X_a)+J_2$, generators $m_1,\ldots,m_g$
of $N$, and $c_{ija}\in O$ such that the $g^2$ rows
$(\bar c_{ija})_a$ are independent and
\[
             \sum_j\Big(\sum_a c_{ija}X_a\Big)m_j\in J_2N
             \qquad(1\le i\le g).
\]
We treat $g>0$. The ideal $\mathfrak m_B+N$ is maximal in $A$
with residue $k_B$. All maximal ideals of $A$ lie over
$\mathfrak m_B$, and $N$ consists there of square-zero elements
modulo $\mathfrak m_B A$, so this maximal ideal is unique.
Since $A$ is finite over $B$, cofinality of its two adic
topologies gives completeness. In $A$, $I_N$ lies in
$\mathfrak m_A^2$. Thus the indicated $X_a,m_j$ span the
relative cotangent space of $A/O$. Nakayama generation on each
successive adic quotient and completeness give a surjection
\[
  F:P=O[[x_1,\ldots,x_{D_0},y_1,\ldots,y_g]]\longrightarrow A,
       \qquad x_a\mapsto X_a,\quad y_j\mapsto m_j .
\]

Write the displayed error as $\sum_j w_{ij}m_j$ with $w_{ij}\in J_2$.
Lift each $w_{ij}$ to $v_{ij}\in(\varpi)+\mathfrak m_P^2$;
this is possible since $F$ is surjective and the image of $J_2$ in $A$
lies in $(\varpi)+\mathfrak m_A^2$.
Set $z_{ij}=\sum_a c_{ija}x_a-v_{ij}$, with images
$\zeta_{ij}=\sum_a c_{ija}X_a-w_{ij}\in B$.
Then $F(zy)=0$. The images of the $z_{ij}$ modulo
$(\varpi,y,\mathfrak m_P^2)$ are independent.
After choosing $D_0-g^2$ complementary $x$-linear forms, the
$z_{ij}$, these forms, and the $y_j$ are formal coordinates
over $O$. Indeed their linear change in the relative cotangent
space is invertible, so successive correction constructs an inverse
in the complete power series algebra (the uniformizer is fixed).
Consequently REX6's generic incidence calculation applies with
the other coordinates in the complete regular base. It says
\[
 H=\sqrt{(zy,\det z)}\text{ is prime},\qquad
              \dim(P/H)=\dim P-g=1+D_0 .
\]
In that calculation the charts with $y_j\ne0$ eliminate a
column of $z$, and a singular matrix at $y=0$ is approached
along a nonzero vector in its kernel. The homogeneous completed
quotient remains a domain by the lowest-degree argument of REX6.

Let $Z,Q\subset B$ be primes with $\dim(B/Z)\ge1+D_0$,
$I_N\not\subseteq Z$, and either $I_N\not\subseteq Q$ or
$\operatorname{Ann}_B N\subseteq Q$. Choose primes $Z',Q'$
of $A$ over them by integrality. Some $m_j$ is outside $Z'$.
Otherwise $N\subset Z'$ gives $I_N\subset Z$. The equations
$zy=0$ over $A/Z'$ then force $\det z=0$ there, so
$H\subset F^{-1}Z'$. The quotient by $F^{-1}Z'$ has dimension
$\dim(B/Z)$ by the integral inclusion $B/Z\hookrightarrow A/Z'$.
Strict containment above $H$ lowers the dimension by at least one
in the finite-dimensional local domain $P/H$. Hence
$F^{-1}Z'=H$.

Multiplying the equation $\zeta m=0$ by the adjugate gives
$\det\zeta\in\operatorname{Ann}_B N$. It therefore belongs to
$Q$ in the second alternative. In the first, it annihilates
$I_N$ since it annihilates $N$, and primality together with
$I_N\not\subseteq Q$ again puts it in $Q$. Thus
$F^{-1}Q'$ contains $H=F^{-1}Z'$. Surjectivity and contraction
give $Z\subseteq Q$. If $g=0$, then $N=0$ and the condition
$I_N\not\subseteq Z$ is impossible. This establishes the stated
implication for all $g$ and isolates the precise missing data at
a vertex: the statements constructing this presentation from
its Selmer conditions and comparing its two ideals with the
discriminant and thick tests are additional mathematical inputs.
% END R70P4

% BEGIN R70P5
\paragraph{R70P5. From the local presentation to the finite-chain theorem.}
Fix any ordinary point and completion $C_i$ of R70P1. Write
$F_i=K_i((t))$. Reduction at the center identifies the residue
field of $C_i$ with $F_i$: the point image contains $K_i$ and $t$
and contains $K_i[[t]]$ by the complete evaluation of $O[[s]]$,
since $O\to K_i$ is onto. The local coefficient DVR is the
completion of $O[[s]]_{(\varpi_O)}$, with residue field $F_i$
and map to $C_i$ induced by the original section. The evaluation
and completion argument of CRW2 uses only the two normalization
conditions, so applies here.

Retain the source's local universal matrix data $U$ at this point
and the source's Selmer presentation $\mathrm{Pres}(D_0)$ in
the same completion, with its labeled $X_a,m_j,c_{ija}$. The
presentation is exactly the maximal-ideal bound, independent
coefficient rows, generators and relations of R70P4. The finite
algebra $A$ and complementary submodule $N$ in that argument
are $U.AA$ and $U.N$, over $B=C_i$. We verify that this suffices
for the local comparison at every allowed block size.

If $a=n$, the data give $N=0$. Completeness and the cotangent
generators give $O_i[[x_1,\ldots,x_{D_0}]]\twoheadrightarrow C_i$,
with $O_i$ the local coefficient DVR just described. A prime
$Z$ with $\dim(C_i/Z)\ge1+D_0$ pulls back to zero in this
regular local domain of dimension $1+D_0$. Thus $Z=0$ and
belongs below every $Q$.

For $0<a<n$ a discriminant prime $Z$ cannot contain $I_N$.
Indeed choose a prime of the integral algebra $A$ over $Z$.
If $I_N\subseteq Z$, the square of every element of $N$ maps
to zero in the residue field upstairs, hence so does $N$.
The matrices of the universal representation specialize to
block diagonals because their off-diagonal entries belong to $N$.
Their span has dimension at most $a^2+(n-a)^2<n^2$, so every
$n^2$-row trace Gram determinant vanishes. The determinant
identity in $U$ identifies these traces with the coefficients
of the given determinant, contradicting the discriminant witness.
Apply this argument also to a discriminant target $Q$. For
a thick target the original $U.\mathrm{thick}$ gives
$\operatorname{Ann}_{C_i} N\subseteq Q$. All the hypotheses
of R70P4 now hold on these exact $Z,Q$, and give $Z\subseteq Q$.

Consequently the two local suppliers, the universal matrix data
and its Selmer presentation at every normalized point, suffice
for the whole $\mathrm{Core}$ and hence the whole $\mathrm{Raw}$
by R70P3. The presentation must be produced from the original
$B1,B2$ and auxiliary eigenvalues. The source passes to $K1,K2,K3$
at the limiting group, then to their Laurent-field versions, before
constructing the tangent generators and their quadratic relations.
R70P4 does not assume the desired prime containment as part of
those input relations. The existing full-block, zero-generator
construction of CRW2--18 supplies its selected case. For general
vertices the two local suppliers remain explicit in this reduction.
% END R70P5

% BEGIN R70P6
\paragraph{R70P6. Full-block vertices in arbitrary rank and chain length.}
At a vertex with $a=n$ the module of odd matrices is zero over
$A=K[[t]]$, its truncations and $F=K((t))$. The package's
independence modulo odd coboundaries implies $g=0$ for
\emph{every} package at that vertex: all its putative generators
are zero, so any index would give the relation with coefficient
one at that index. The coefficient of the first Selmer problem is
then exactly $M_n$, with its prescribed polarized boundaries.
We show that the two local suppliers in R70P5 are automatic here
for every positive $n$.

The hypothesis on the span of the point gives the full matrix
algebra over $F$ after inverting $t$. The absolute irreducibility,
henselian determinant representation, polarization and ordinary
relation argument of CRW3--6 therefore gives the actual $U$ with
$A_U=C_i$ and $N=0$. Those steps use the full $n$-by-$n$ matrix
span, the trace pairing on matrix units and the oddness of $\ell$.
The residue representation, ordered roots, ordinary point and local
DVR to which they apply have been identified in R70P1 and R70P5.

For the Selmer comparison fix any truncation. The common limit
of the Taylor--Wiles groups has a finitely generated profinite
cover. For its fixed finite action on $M_n(A/t^m)$ there are
only finitely many continuous cocycles, since each is determined
on the finite set of topological generators by its finite affine
action and values. On an ultrafilter-large set of levels the
projection, kernel and local convergence conditions of the limit
identify all their descent tests and ordinary conditions, as in
CRW12. Their boundaries are the images of the same finite matrices.
Intersect that set with the eventual $B1$ inequality supplied at
this vertex, giving the same cardinal bound at the limit and
the same single exponent before every truncation (CRW13).
The finite DVR system of CRW14 tests whether a linear combination
of cocycles is a polarized coboundary, including the skew primitive
variable. Smith normal form bounds the loss in every truncated
independence test by a constant. As $m$ grows the $B1$ bound
therefore limits the number of independent classes to $D_0$.
Clearing denominators as in CRW15 gives this same bound over $F$.
For $K2$ the sole coefficient row can have coefficient one and
zero lifting cocycle, since it has zero columns. For $K3$ the
odd Selmer module is zero. The same choices apply over $F$.
Thus no general finite-level-to-limit Selmer theorem is required
at a full-block vertex.

The cotangent-to-cocycle comparison in CRW7--11 applies to these
$n$-by-$n$ matrices and their full span. In particular elementary
linear algebra suffices here: a relative derivation whose cocycle
is a boundary kills the characteristic coefficients, since an
infinitesimal conjugation fixes them. The local ordinary characters
and the auxiliary simple-root coefficients are then fixed by
their ordered first-order relations, as checked in CRW10--11.
The resulting injection on the relative cotangent dual and the
preceding $F$-bound give at most $D_0$ generators. Lift a basis
and pad with zeros as in CRW16. With $N=0$ all odd generators,
quadratic relations and coefficient-row conditions are empty.
This proves the full presentation at this vertex. R70P5 gives
its local comparison.

In particular every chain consisting entirely of full-block points
satisfies the original conclusion. In a mixed chain the comparison
is now needed only at the vertices with $0<a<n$. At those vertices
the limiting implication from the given $B1,B2$ and eigenvalues to
the three integral Selmer tests, their passage to Laurent
coefficients, and the construction of the presentation from those
tests are still required if one follows the source route. R70P1--3
use the resulting prime comparison at such vertices and R70P4--5
prove it from the indicated presentation and universal matrix data.
% END R70P6
